\section{Aplikacja testowa}
\label{sec:aplikacja}

\subsection{Wymagania}
\label{sec:wymagania}
%<*wymagania>

Wymagania funkcjonalne:
\begin{enumerate}[label=F\arabic*.]
  \item Generowanie tego samego syntetycznego zbioru danych o~zadanym rozmiarze
  (liczba biletów) i~ziarnie dla każdej bazy oraz potwierdzenie identyczności
  danych sumą kontrolną.
  \item Import hurtowy danych do PostgreSQL, MariaDB, MongoDB i~Valkey.
  \item Wykonanie każdego z~27 scenariuszy na każdej bazie z~zadaną liczbą
  powtórzeń, w~wariancie bez indeksów drugorzędnych i~z~nimi.
  \item Zapis planów zapytań dla każdego scenariusza i~wariantu indeksów.
  \item Test współbieżności: $N$ klientów rezerwujących te same miejsca, wiele
  strategii blokowania, kontrola braku podwójnej sprzedaży.
  \item Zapis surowych wyników do CSV i~wygenerowanie z~nich tabel, wykresów
  oraz oceny hipotez bez ręcznego przepisywania liczb.
  \item Uruchomienie całości jednym poleceniem.
\end{enumerate}

Wymagania niefunkcjonalne:
\begin{enumerate}[label=N\arabic*.]
  \item Powtarzalność: deterministyczne dane i~parametry, wersje obrazów
  i~konfiguracja baz zapisane w~repozytorium.
  \item Rzetelność pomiaru: czas mierzony w~serwisie bazy, bez narzutu sieci
  między CLI a~serwisem; rozgrzewka przed pomiarem; fikstury poza pomiarem.
  \item Obserwowalność: każdy wynik powiązany z~identyfikatorem śladu
  OpenTelemetry, metryki w~Prometheusie.
  \item Przenośność: uruchomienie na macOS i~Linuksie (amd64, arm64) w~Dockerze.
  \item Jakość: testy automatyczne na prawdziwych bazach w~CI przy każdej zmianie.
\end{enumerate}

%</wymagania>
\subsection{Technologie i narzędzia}
\label{sec:technologie}

\begin{table}[H]
\centering
\small
\caption{Wykorzystane technologie.}
%<*technologie>
\begin{tabular}{>{\raggedright\arraybackslash}p{3.6cm}>{\raggedright\arraybackslash}p{11cm}}
\toprule
obszar & technologie \\
\midrule
język i~komunikacja & Go 1.26, gRPC i~Protocol Buffers (buf), cobra (CLI) \\
sterowniki baz & pgx v5 (PostgreSQL), go-sql-driver/mysql (MariaDB), mongo-driver v2, valkey-go \\
obserwowalność & OpenTelemetry SDK, otelgrpc, otelpgx, otelsql, otelmongo, valkeyotel, OpenTelemetry Collector, Jaeger, Prometheus, Grafana \\
środowisko & Docker, Docker Compose, obrazy distroless, GHCR \\
testy & testing (Go), testcontainers-go, golangci-lint, ansible-lint, shellcheck \\
CI/CD i~wdrożenie & GitHub Actions, GoReleaser, Ansible, skrypt \texttt{ticketbench.sh} \\
raport & LaTeX (XeLaTeX, tectonic), pgfplots, siunitx; tabele generowane przez \texttt{tbench report} \\
\bottomrule
\end{tabular}
%</technologie>
\end{table}

\subsection{Architektura}
\label{sec:architektura}

Aplikacja \emph{ticketbench} jest napisana w~Go (1.26) i~podzielona na
mikroserwisy komunikujące się przez gRPC (rys.~\ref{fig:arch}). Kontrakty
(\texttt{StoreService}, \texttt{RunnerService}), model danych, generator,
definicje scenariuszy i~konfiguracja OpenTelemetry tworzą wspólny moduł
\texttt{ticketbench-contracts}.

\begin{figure}[H]
\centering
\begin{tikzpicture}[
  box/.style={draw, rounded corners=3pt, align=center, font=\footnotesize, minimum height=0.85cm, inner sep=4pt},
  svc/.style={box, fill=blue!6},
  db/.style={box, fill=gray!10},
  obs/.style={box, fill=orange!10},
  arr/.style={-{Stealth[length=2mm]}},
  otlp/.style={-{Stealth[length=2mm]}, dashed, gray}]
  \node[box, fill=green!8] (cli) {\texttt{tbench} CLI\\root span};
  \node[svc, right=1.3cm of cli] (run) {bench-runner\\plan, wyniki};
  \node[svc, right=1.6cm of run, yshift=1.9cm] (spg) {store-postgres};
  \node[svc, below=0.35cm of spg] (smy) {store-mariadb};
  \node[svc, below=0.35cm of smy] (smo) {store-mongo};
  \node[svc, below=0.35cm of smo] (sva) {store-valkey};
  \node[db, right=1.1cm of spg] (pg) {PostgreSQL 17};
  \node[db, right=1.1cm of smy] (my) {MariaDB 11.8};
  \node[db, right=1.1cm of smo] (mo) {MongoDB 8 (rs0)};
  \node[db, right=1.1cm of sva] (va) {Valkey 8.1};
  \node[obs, below=1.3cm of run] (col) {OTel Collector};
  \node[obs, left=0.7cm of col] (jae) {Jaeger};
  \node[obs, below=0.5cm of col, xshift=-1.2cm] (pro) {Prometheus};
  \node[obs, right=0.5cm of pro] (gra) {Grafana};
  \draw[arr] (cli) -- node[above, font=\scriptsize]{gRPC} node[below, font=\scriptsize]{traceparent} (run);
  \foreach \s in {spg,smy,smo,sva} \draw[arr] (run.east) -- (\s.west);
  \draw[arr] (spg) -- (pg); \draw[arr] (smy) -- (my); \draw[arr] (smo) -- (mo); \draw[arr] (sva) -- (va);
  \draw[otlp] (run) -- (col); \draw[otlp] (sva.south) |- (col.east);
  \draw[otlp] (col) -- (jae); \draw[otlp] (col) -- (pro); \draw[otlp] (pro) -- (gra);
\end{tikzpicture}
\caption{Architektura: linie ciągłe to wywołania gRPC i~połączenia z~bazą, przerywane to eksport OTLP.}
\label{fig:arch}
\end{figure}

%<*dzialanie>
\begin{description}
  \item[store-*] Jeden serwis na bazę, wspólny interfejs \texttt{Engine} w~Go.
  Serwis generuje dane z~ziarna, ładuje je hurtowo (PostgreSQL: \texttt{COPY}
  w~formacie binarnym, MariaDB: wielowierszowe \texttt{INSERT} po 1000 wierszy,
  MongoDB: \texttt{insertMany} po 5000 dokumentów, Valkey: potoki poleceń),
  tworzy i~usuwa indeksy, wykonuje scenariusze, zwraca plany zapytań
  i~przeprowadza test współbieżności.
  \item[bench-runner] Przyjmuje plan (listę kroków), wykonuje go po kolei
  i~strumieniuje wiersze wyników do klienta.
  \item[tbench] Klient CLI: buduje plan z~pliku \texttt{matrix.yaml}, zapisuje
  wyniki do CSV na bieżąco, generuje tabele i~dane wykresów do tego
  sprawozdania (\texttt{tbench report}).
\end{description}
%</dzialanie>

\subsection{Śledzenie i trace id}
\label{sec:sledzenie}
%<*sledzenie>

CLI rozpoczyna span główny dla każdego segmentu (jeden rozmiar zbioru na
jednej bazie). Kontekst śladu trafia do runnera jako nagłówek W3C
\texttt{traceparent} w~metadanych gRPC~\cite{w3c}. Runner tworzy spany kroków
i~wywołań \texttt{Execute}, przekazuje kontekst dalej do serwisu bazy, a~ten
do instrumentacji sterownika (otelpgx, otelsql, otelmongo, valkeyotel). Każde
polecenie SQL, operacja MongoDB i~polecenie Valkey jest więc spanem w~tym
samym śladzie. Logi (\texttt{slog}, JSON) zawierają pola \texttt{trace\_id}
i~\texttt{span\_id}. Identyfikator śladu jest zapisany w~każdym wierszu
plików CSV z~wynikami, więc z~dowolnego pomiaru można przejść do jego śladu
w~Jaegerze. Test jednostkowy runnera sprawdza, że nagłówek
\texttt{traceparent} z~identyfikatorem śladu klienta dociera do serwisu bazy.

Przy imporcie danych serwis oznacza kontekst jako niepróbkowany
(\texttt{telemetry.Quiet}): identyfikator śladu zostaje, ale tysiące spanów
pojedynczych partii \texttt{INSERT} nie są zapisywane. W~teście
współbieżności śledzimy 1\% prób rezerwacji. Metryki (histogramy czasu
scenariuszy i~prób rezerwacji, licznik konfliktów) trafiają przez kolektor do
Prometheusa i~na pulpit Grafany.

%</sledzenie>
\subsection{Automatyzacja}
\label{sec:testy}

%<*automatyzacja>
\begin{itemize}
  \item \textbf{Docker Compose} uruchamia 4 bazy, 4 serwisy, runner, kolektor
  OTel, Jaeger, Prometheus i~Grafanę. Obrazy serwisów są oparte na
  \texttt{distroless/static} (ok.~30\,MB) i~mają kontrolę zdrowia przez
  gRPC Health.
  \item \textbf{Jeden program} (\texttt{ticketbench.sh}) pobiera wszystkie
  serwisy z~GitHuba (repozytoria, obrazy, binarkę \texttt{tbench}), uruchamia
  środowisko i~całą macierz testów, generuje tabele i~wykresy, a~na koniec
  kompiluje to sprawozdanie i~prezentację. Instrukcja: runbook
  \texttt{docs/runbook.html}.
  \item \textbf{Testy automatyczne}: testy jednostkowe generatora (determinizm,
  spójność rezerwacji i~biletów), scenariuszy, generatora raportu i~CLI,
  a~także test zgodności (\texttt{enginetest}) uruchamiany na prawdziwych
  bazach przez testcontainers-go. Test zgodności porównuje wyniki zapytań
  z~wyrocznią liczoną bezpośrednio z~generatora, wykonuje wszystkie
  scenariusze w~obu wariantach indeksów, plany zapytań i~test współbieżności
  każdą strategią.
  \item \textbf{GitHub Actions}: lint (golangci-lint, buf, ansible-lint),
  testy z~detektorem wyścigów, test zgodności jako macierz czterech baz,
  benchmarki Go, binarki dla linux/amd64, linux/arm64 i~darwin/arm64,
  wieloarchitekturowe obrazy w~GHCR, wydania przez GoReleaser oraz test
  end-to-end, w~którym Ansible stawia środowisko na maszynie CI i~uruchamia
  macierz \texttt{smoke}. Sprawozdanie również kompiluje się w~CI.
  \item \textbf{Wariant bezpieczeństwa}: nakładka
  \texttt{docker-compose.security.yml} (certyfikaty z~własnego CA, konta
  aplikacyjne) i~polecenie \texttt{./ticketbench.sh security}.
  \item \textbf{Ansible}: role \texttt{docker} (instalacja silnika na
  Debianie/Ubuntu), \texttt{ticketbench} (wdrożenie i~start stosu),
  \texttt{bench} (uruchomienie macierzy i~pobranie wyników)
  i~\texttt{teardown}.
\end{itemize}
%</automatyzacja>
